\documentclass{article}
\usepackage{amsmath}
\usepackage{amssymb}
\usepackage{hyperref}
\usepackage{float}
\usepackage[swedish]{babel}
\usepackage{bm}
\newcommand{\indep}{\perp \!\!\! \perp}
\newcommand{\notindep}{\not\!\perp\!\!\!\perp}
\newcommand{\var}{\text{Var}}
\newcommand{\cov}{\text{Cov}}
\newcommand{\corr}{\text{Corr}}
\newcommand{\set}[1]{\left\{#1\right\}}
\newcommand{\an}{\text{an}}
\newcommand{\pa}{\text{pa}}

\title{StatDat Projekt}
\author{Christoffer Teorell}
\date{September 2026}

\begin{document}

\maketitle

\section{Del A. Ett land, två grupper}

\subsection{A1. Förhandskontroll och förregistrering}

\subsection{A2. Frekventistisk utgångspunkt}

Vi har $X_j|\Theta_j=\theta_j \sim \text{Bin}(n_j,\theta_j)$ och samma för k. Vi beräknade att

\[
\begin{cases}
    \hat \theta_{j,\text{obs}}=0.194\\
    n_j = 325\\
    \hat \theta_{k,\text{obs}} = 0.24\\
    n_k = 546
\end{cases}
\]

\subsubsection{ML-Skattning av delta}

Vi antar $X_j|\Theta_j=\theta_j\indep X_k|\Theta_k=\theta_k$

Vi vill visa att vi har följande ML-skattning

\[
\hat \delta =\hat \theta_j-\hat \theta_{k}= \frac{X_j}{n_j}-\frac{X_k}{n_k}
\]

Vi börjar med att visa ML-skattningen för p för en Bin($n,p$) fördelning. Vi har

\[
\begin{aligned}
&L(x|p)=\begin{pmatrix}
    n\\x
\end{pmatrix}p^x(1-p)^{n-x}\\
&\implies \log L(x|p)=\log \begin{pmatrix}
    n\\x
\end{pmatrix}+x\log p +(n-x)\log (1-p)\\
&\implies \frac{\partial }{\partial p}\log L(x|p)=\frac{x}{p}-\frac{n-x}{1-p}=0\\
&\implies \frac{x}{p}=\frac{n-x}{1-p}\implies x-px=np-px \implies p=\frac{x}{n}
\end{aligned}
\]

Eftersom vi har oberoende kan vi skriva vår simultana  täthetsfunktionen som

\[
f_{X_j,X_k|\Theta_j,\Theta_k}(x_j,x_k|\theta_j,\theta_k)=f_{X_j|\Theta_j}(x_j|\theta_j)f_{X_k|\Theta_k}(x_k|\theta_k)
\]

därmed får vi följande log-likelihood funktion

\[
\log L(x_j,x_k|\theta_j,\theta_k)=\log f_{X_j|\Theta_j}(x_j|\theta_j)+\log f_{X_k|\Theta_k}(x_k|\theta_k)
\]

Vilket enligt ovan ger ML-skattningen för $[\theta_j,\theta_k]^T$

\[
\begin{cases}
    \hat \theta_{j}=\frac{X_j}{n_j}\\
    \hat \theta_{k}=\frac{X_k}{n_k}
\end{cases}
\]

Vi kan nu använda invariansegenskapen av ML-skattningar. Låt $\tau: \mathbb{R}^2\rightarrow \mathbb{R}^2$ vara definierad av

\[
\tau(a,b)=(a-b,b)
\]

Vi noterar att $\tau $ är bijektiv med invers $\tau^{-1}(a,b)=(a+b,b)$. Då ger invariansegenskapen att ML-skattningen av $\tau([\theta_j,\theta_k])$ är 

\[
\tau(\hat \theta_{j},\hat \theta_k)=(\hat \theta_j-\hat \theta_k,\hat \theta_k)=(\hat \delta,\hat \theta_k)
\]
Vi kan läsa av första raden och ser att ML-skattningen av $\theta_j-\theta_k$ ges av $\hat \delta =\hat \theta_j-\hat \theta_k$

Skattningen säger följande: skillnaden mellan andelen positiva svar och andelen negativa svar approximeras av 

\subsubsection{Väntevärdesriktighet och Cramér-Rao}

Vi har skattningen

\[
\hat \delta =\frac{1}{n_j}X_j-\frac{1}{n_k}X_k
\]

Väntevärdet av skattningen ges av

\[
\begin{aligned}
    E[\hat \delta]&=\frac{1}{n_j}E[X_j  ]-\frac{1}{n_k}E[X_k]\\
    &=\frac{1}{n_j}n_j\theta_j-\frac{1}{n_k}n_k\theta_k\\
    &=\theta_j-\theta_k
\end{aligned}
\]

Vi hittar nu variansen

\[
\begin{aligned}
    \var [\hat \delta ]&=\frac{1}{n_j^2}\var[X_j]+\frac{1}{n_k^2}\var[X_k]\\
    &=\frac{1}{n_j^2}\big(n_j\theta_j(1-\theta_j)\big)+\frac{1}{n_k^2}\big(n_k\theta_k(1-\theta_k)\big)\\
    &=\frac{\theta_j(1-\theta_j)}{n_j}+\frac{\theta_k(1-\theta_k)}{n_k}
\end{aligned}
\]


Cramér Rao säger att

\[
\var[W[X]]\geq \nabla g^TI^{-1}(\theta_j,\theta_k)\nabla g
\]
Vi har

\[
g(\theta_j,\theta_k)=\delta =\theta_j-\theta_k\implies \nabla g=(1,-1)
\]
\[
W[X]=\hat \delta =\frac{X_j}{n_j}+\frac{X_k}{n_k}
\]

Fischer informationen ges av följande matris

\[
I(\theta_j,\theta_k)=\begin{pmatrix}
    E\Bigg[\Big(\frac{\partial }{\partial \theta_j}\log L\Big)^2\Bigg]& E\Bigg[\frac{\partial }{\partial \theta_j}(\log L)\frac{\partial }{\partial \theta_k}(\log L)\Bigg]\\ E\Bigg[\frac{\partial }{\partial \theta_k}(\log L)\frac{\partial }{\partial \theta_j}(\log L)\Bigg] & E\Bigg[\Big(\frac{\partial }{\partial \theta_k}\log L\Big)^2\Bigg]
\end{pmatrix}
\]

Vi kan även skriva det i termer av andraderivator och får då

\[I(\theta_j,\theta_k)=-
\begin{pmatrix}
    E\Bigg[\frac{\partial^2}{\partial \theta_j^2}\log L\Bigg]&E\Bigg[\frac{\partial^2}{\partial \theta_j\partial \theta_k}\log L\Bigg] \\
    E\Bigg[\frac{\partial^2}{\partial \theta_k\partial \theta_j}\log L\Bigg]&E\Bigg[\frac{\partial^2}{\partial \theta_k^2}\log L\Bigg]
\end{pmatrix}
\]

Vi har

\[
\begin{aligned}
     \frac{\partial }{\partial \theta_j}\log L&=\frac{X_j}{\theta_j}-\frac{n_j-X_j}{1-\theta_j}\\
     \implies \frac{\partial ^2}{\partial \theta_j^2}\log L &=-\frac{X_j}{\theta_j^2}-\frac{n_j-X_j}{(1-\theta_j)^2}
\end{aligned}
\]

Alltså får vi följande väntevärde

\[
E\Bigg[\frac{\partial^2}{\partial \theta_j^2}\log L\Bigg]=-\frac{1}{\theta_j^2}E[X_j]-\frac{n_j}{(1-\theta_j)^2}+\frac{1}{(1-\theta_j)^2}E[X_j]
\]

För fixt $\theta_j$ har vi att $E[X_j|\Theta_j=\theta_j]=n_j\theta_j$. Insättning ger


\[
\begin{aligned}
    E\Bigg[\frac{\partial^2}{\partial \theta_j^2}\log L\Bigg]&=-\frac{n_j\theta_j}{\theta_j^2}-\frac{n_j}{(1-\theta_j)^2    }+\frac{n_j\theta_j}{(1-\theta_j)^2}\\
    &=-\frac{n_j}{\theta_j  }-\frac{n_j(1-\theta_j)}{(1-\theta_j)^2}\\
    &=-\Big(\frac{n_j}{\theta_j}+\frac{n_j}{(1-\theta_j)}\Big)=\frac{-n_j}{\theta_j(1-\theta_j)}
\end{aligned}
\]
Detta ger även att

\[
 E\Bigg[\frac{\partial^2}{\partial \theta_k^2}\log L\Bigg]=\frac{-n_k}{\theta_k(1-\theta_k)}
\]

Vidare blir blandtermerna lika med 0 eftersom $\frac{\partial}{\partial\theta_j}\log L$ inte innehåller $\theta_k$ och vice versa

Alltså får vi följande Fisher-informationsmatris

\[
I(\theta_j,\theta_k)=-\begin{pmatrix}
    \frac{-n_j}{\theta_j(1-\theta_j)}&0\\0&\frac{-n_k}{\theta_k(1-\theta_k)}
\end{pmatrix}
\]

Alltså har vi följande invers:

\[
I^{-1}(\theta_j,\theta_k)=\begin{pmatrix}
    \frac{\theta_j(1-\theta_j)}{n_j}&0\\0&\frac{\theta_k(1-\theta_k)}{n_k}
\end{pmatrix}
\]

Alltså har vi att den undre gränsen från Cramér Raos sats ges av

\[
\nabla g^TI^{-1}(\theta_j,\theta_k)\nabla g=\frac{\theta_j(1-\theta_j)}{n_j}+\frac{\theta_k(1-\theta_k)}{n_k}
\]

Vi ser nu att detta är precis lika med $\var[\delta]$. Därmed uppfyller vi den lägre gränsen för Cramér Rao


\subsection{A3. Konjugerade aposteriorifördelningar}

\subsubsection{Härledning av aposteriorifördelning}

Vi har följande uppställning

\[
\begin{cases}
    X|\Theta=\theta\sim \text{Bin}(n,\theta)\\
    \Theta\sim \text{Beta}(\alpha,\beta)
\end{cases}
\]

Vi vill hitta aposteriorifördelningen som ges av

\[
f_{\Theta|X}(\theta|x)=\frac{f_{X|\Theta}(x|\theta)f_\Theta(\theta)}{f_X(x)}
\]

Vi vet att 

\[
f_X(x)=\int_\mathbb{\Omega}f_{X|\Theta}(x|\theta)f_{\Theta}(\theta)d\theta \in \mathbb{R}
\]

Eftersom $f_X(x)\in \mathbb{R}$ kan vi ignorera den och betrakta följande proportionalitetspåstående istället

\[
f_{\Theta|X}(\theta|x)\propto f_{X|\Theta}f_{\Theta}(\theta)
\]

Vi vet att 

\[
\begin{cases}
    f_{X|\Theta}(x|\theta)=\begin{pmatrix}
        n\\x
    \end{pmatrix}\theta^x(1-\theta)^{n-x}\\
    f_{\Theta}(\theta)=\frac{\Gamma(\alpha+\beta)}{\Gamma(\alpha)\Gamma(\beta)}\theta^{\alpha-1}(1-\theta)^{\beta-1}
\end{cases}
\]

Då får vi att

\[
f_{\Theta|X}(\theta|x)\propto \begin{pmatrix}
    n\\x    
\end{pmatrix} \frac{\Gamma(\alpha+\beta)}{\Gamma(\alpha)\Gamma(\beta)}\theta^{(\alpha+x)-1  }(1-\theta)^{(n+\beta-x)-1}
\]

Vi kan ignorera konstanten och identifiera detta med en Beta$(\alpha+x,n+\beta-x)$ fördelning. Vi kan nu sätta $\alpha=\beta=0$ och får

\[
\Theta|X\sim \text{Beta}(x,n-x)
\]

Därmed har vi följande aposteriorifördelningar för vårt $\Theta_j,\Theta_k$

\[
\begin{cases}
    \Theta_j|X_j\sim \text{Beta}(x_j,n_j-x_j)\\
    \Theta_k|X_k\sim \text{Beta}(x_k,n_k-x_k)
\end{cases}
\]


\subsubsection{Simulering av Delta}

Free, svara frpgan om varför

\subsubsection{Bayesianskt hypotestest}

Vi formar hypotesen

\[
H_0:\theta_j\leq \theta_k \text{ mot } H_1:\theta_j>\theta_k
\]
med förlustfunktionen

\[
L(
\]

ssss




\end{document}


